% !TeX root = ../mub6_karlsson_ieee.tex
% Theorem T2: Third MUB on the Dita lambda-circle.
% Include from main paper or compile standalone.

\section{Third MUB on the Dita \texorpdfstring{$\lambda$}{lambda}-circle (T2)}
\label{sec:proof-t2}

\begin{theorem}[Third MUB on Dita slice, T2 (HP-supported)]
\label{thm:T2}
Let $\theta_D=\arccos(1/\sqrt3)$ and $\phi=\pi/4$.
Write $H(\lambda):=H(\theta_D,\phi,\lambda)$.
At each tested $\lambda$ on the Dita circle ($126/126$ periodicity samples on $[0,2\pi]$ and $628/628$ dense probes), the pair $\{\mathbf{I},H(\lambda)\}$ admits a third mutually unbiased basis:
there exist six unit vectors forming an orthonormal basis, each unbiased to
$\mathbf{I}$ and to every column of $H(\lambda)/\sqrt6$, obtained by certified pool extraction.
By $2\pi$-periodicity (Theorem~\ref{thm:T1}, L2) and the informal continuity argument below (Route~B), the same conclusion is inferred for all $\lambda\in\mathbb{R}$; that extension is not a closed-form or purely analytic proof.
\end{theorem}

\begin{proof}[Constructive HP certificate]
Fix $\lambda$.
Solve the per-$H$ MU-pool polynomial system (10 equations, 10 variables) at $H(\lambda)$
via certified homotopy continuation (\texttt{generate\_candidate\_pool\_fresh}).
Let $\mathcal{P}(\lambda)$ be the deduplicated pool of MU vectors.

Define the orthogonality graph $G(\lambda)$ on $\mathcal{P}(\lambda)$ with edges
when $|\langle v,w\rangle|<\varepsilon_{\mathrm{ortho}}$.
By \texttt{extract\_third\_mub\_basis} (\texttt{src/dita\_third\_mub\_construction.jl}),
if $G(\lambda)$ contains a $6$-clique $C$, the columns of $B_3=[v_{c_1},\ldots,v_{c_6}]$
form the required third ONB.

\textbf{Anchor.} At $\lambda_0=0.4$, Brierley--Weigert established third MUB existence at Dita-type parameters (Section~\ref{sec:prior}); our pipeline recovers \texttt{max\_clique}$=6$.

\textbf{Parametric persistence.} Theorem~\ref{thm:T1} (L2) gives
$H(\lambda+2\pi)=H(\lambda)$. High-precision verification
(\texttt{verify\_dita\_construction.jl}, \texttt{lambda\_periodicity\_dita.jl}):
\begin{itemize}
    \item $126/126$ samples on $[0,2\pi]$ with clique $\ge6$;
    \item $628/628$ dense $\lambda$ probes with \texttt{found\_fourth=false};
    \item HP re-verification (256--400 bit) of clique orthogonality at spot checks.
\end{itemize}
The explicit construction $\lambda\mapsto B_3(\lambda)$ is computed independently at each
$\lambda$; CHM-inequivalence of distinct $H(\lambda)$ (Lemma~\ref{lem:L4}) shows this is
not a single-matrix artifact.

\textbf{Informal extension (Route B).}
On each open interval between $\lambda$-values where pool completeness holds,
the clique template varies continuously when roots are certified.
Combined with $2\pi$-periodicity and the full-circle HP audit ($126/126$, $628/628$),
this suggests third MUB existence for all $\lambda\in\mathbb{R}$, but the continuity step is not a rigorous analytic proof.
\end{proof}

\begin{remark}
A fully closed-form symbolic formula for $B_3(\lambda)$ from Brierley--Weigert's
Gr\"obner construction remains future work; T2 is supported by certified constructive
extraction at each tested $\lambda$, plus the full-circle HP audit ($126/126$, $628/628$); a closed-form or analytic proof for all $\lambda\in\mathbb{R}$ remains open.
\end{remark}
